由引理~\ref{ms:lemma-3-2} 的秩一行列式恒等式，
\[
 \det D_{\theta u}
 =\theta^n\det D_u+(1-\theta)\theta^{n-1}C(D_u):B_{\rm fib}.
\]
令 $I=\int_{\Omega_z}(-u)C(D_u):B_{\rm fib}\,\mathrm d\zeta\ge0$。对该多项式关于 $\theta$ 积分，得到 $\mathcal F=nX_0/(n+1)-I/[\lambda n(n+1)]$，从而得上界。齐次性给出 $C(D_u):D_u=n\det D_u$，故
\[
 n\lambda X_0-I=\int_{\Omega_z}(-u)C(D_u):H_{\rm fib}(u)\,\mathrm d\zeta
 =\int_{\Omega_z}G_{\rm fib}(u)^*C(D_u)G_{\rm fib}(u)\,\mathrm d\zeta.
\]
在最后一个等式中，编码后的梯度分量为 $G_{{\rm fib},j}(u)=\partial_{j0}u+\mathbf i\partial_{j1}u+\mathbf j\partial_{j2}u$。在这些槽中，$C(D_u)$ 的实形式表示作用于此梯度的四元数二次型。由引理~\ref{ms:lemma-3-3}，它的每个实列的散度为零，而边界项由 $u=0$ 而消失。代入即得恒等式和下界。
